\begin{ex}
Tính: $354200 + 245800 = \dots$
\shortans{600000}
\loigiai{
$354200 + 245800 = 600000$.
}
\end{ex}

\begin{ex}
Tính: $789500 - 345200 = \dots$
\shortans{444300}
\loigiai{
$789500 - 345200 = 444300$.
}
\end{ex}

\begin{ex}
Tính: $1000000 - 350000 = \dots$
\shortans{650000}
\loigiai{
$1000000 - 350000 = 650000$.
}
\end{ex}

\begin{ex}
Tính nhanh: $45000 + 28000 + 55000$.
\shortans{128000}
\loigiai{
$(45000 + 55000) + 28000 = 100000 + 28000 = 128000$.
}
\end{ex}

\begin{ex}
Tính nhanh: $125000 + 340000 + 875000$.
\shortans{1340000}
\loigiai{
$(125000 + 875000) + 340000 = 1000000 + 340000 = 1340000$.
}
\end{ex}

\begin{ex}
Tính giá trị biểu thức: $a + b - c$ với $a = 500000, b = 300000, c = 200000$.
\shortans{600000}
\loigiai{
$500000 + 300000 - 200000 = 600000$.
}
\end{ex}

\begin{ex}
Tìm hai số có tổng là 120 và hiệu là 20. Số lớn là bao nhiêu?
\shortans{70}
\loigiai{
$(120 + 20) : 2 = 70$.
}
\end{ex}

\begin{ex}
Tìm hai số có tổng là 120 và hiệu là 20. Số bé là bao nhiêu?
\shortans{50}
\loigiai{
$(120 - 20) : 2 = 50$.
}
\end{ex}

\begin{ex}
Một hình chữ nhật có nửa chu vi là 60 m, chiều dài hơn chiều rộng 16 m. Chiều dài là bao nhiêu mét?
\shortans{38}
\loigiai{
$(60 + 16) : 2 = 38$ m.
}
\end{ex}

\begin{ex}
Chiều rộng hình chữ nhật trong bài toán trên là bao nhiêu mét?
\shortans{22}
\loigiai{
$(60 - 16) : 2 = 22$ m.
}
\end{ex}

\begin{ex}
Tìm $x$: $x + 350000 = 800000$.
\shortans{450000}
\loigiai{
$x = 800000 - 350000 = 450000$.
}
\end{ex}

\begin{ex}
Tìm $x$: $x - 240000 = 560000$.
\shortans{800000}
\loigiai{
$x = 560000 + 240000 = 800000$.
}
\end{ex}

\begin{ex}
Tìm $x$: $1000000 - x = 420000$.
\shortans{580000}
\loigiai{
$x = 1000000 - 420000 = 580000$.
}
\end{ex}

\begin{ex}
Một trường có 650 học sinh, số học sinh nữ nhiều hơn số học sinh nam là 50 bạn. Hỏi trường có bao nhiêu học sinh nữ?
\shortans{350}
\loigiai{
$(650 + 50) : 2 = 350$ học sinh nữ.
}
\end{ex}

\begin{ex}
Trường có bao nhiêu học sinh nam trong bài toán trên?
\shortans{300}
\loigiai{
$(650 - 50) : 2 = 300$ học sinh nam.
}
\end{ex}

\begin{ex}
Hai kho thóc có 450 tấn thóc, kho thứ nhất có nhiều hơn kho thứ hai 50 tấn thóc. Kho thứ nhất có bao nhiêu tấn thóc?
\shortans{250}
\loigiai{
$(450 + 50) : 2 = 250$ tấn.
}
\end{ex}

\begin{ex}
Kho thứ hai có bao nhiêu tấn thóc trong bài toán trên?
\shortans{200}
\loigiai{
$(450 - 50) : 2 = 200$ tấn.
}
\end{ex}

\begin{ex}
Tính giá trị biểu thức: $1000000 - (450000 + 250000)$.
\shortans{300000}
\loigiai{
$1000000 - 700000 = 300000$.
}
\end{ex}

\begin{ex}
Tính nhanh: $36 + 48 + 64 + 52$.
\shortans{200}
\loigiai{
$(36 + 64) + (48 + 52) = 100 + 100 = 200$.
}
\end{ex}

\begin{ex}
Tính: $999999 - 333333 = \dots$
\shortans{666666}
\loigiai{
$999999 - 333333 = 666666$.
}
\end{ex}

\begin{ex}
Tính: $4000000 + 6000000 = \dots$
\shortans{10000000}
\loigiai{
$4000000 + 6000000 = 10000000$.
}
\end{ex}

\begin{ex}
Tuổi bố và con cộng lại được 50 tuổi, bố hơn con 30 tuổi. Bố bao nhiêu tuổi?
\shortans{40}
\loigiai{
$(50 + 30) : 2 = 40$ tuổi.
}
\end{ex}

\begin{ex}
Con bao nhiêu tuổi trong bài toán trên?
\shortans{10}
\loigiai{
$(50 - 30) : 2 = 10$ tuổi.
}
\end{ex}

\begin{ex}
Tính: $850000 - 350000 + 200000 = \dots$
\shortans{700000}
\loigiai{
$500000 + 200000 = 700000$.
}
\end{ex}

\begin{ex}
Tổng của số chẵn lớn nhất có 3 chữ số (998) và số chẵn bé nhất có 3 chữ số (100) là bao nhiêu?
\shortans{1098}
\loigiai{
$998 + 100 = 1098$.
}
\end{ex}
